\section{Security Analysis}
\label{sec:secannalysis}

\begin{figure}[t]
    \centering
    \includegraphics[width=\columnwidth]{figures/memcpy_code.pdf}
    \caption{Speculative control flow of \texttt{memcpy} assembly under our two mechanisms for secure shared memory}
    \label{fig:memcpyasm}
    \end{figure}
    
Citadel builds upon MI6, which already enforces Strong Microarchitectural Isolation (SMI) through partitioning and flushing
(See Section \ref{sec:hardwareisolation} and \cite{bourgeat2019mi6} for details).
As a result, our security analysis focuses on the novelties brought forward by Citadel,
i.e., shared memory between an enclave and the OS.
This change effectively opens side channels in the shared memory hierarchy through the shared LLC partitions.
Note that any other side channels (e.g., branch predictors) are still closed per inherited MI6's isolation mechanisms: partitioning and flushing. 
Furthermore, we only worry about the confidentiality of the enclave with regard to the OS here (see Section \ref{sec:threatmodel}).
As explained in Section \ref{sec:threatmodel}, the high-level security guarantee we are trying to enforce is Speculative Non-Interference \cite{guarnieri2020spectector} (SNI) between the attacker and the victim program.
To formally reason about SNI and capture these new interactions, we need a different set of arguments than for SMI and to leverage the framework of hardware software contracts for security introduced in Section \ref{sec:hwswcontract}. 
Following this methodology, we first define the security contracts we will consider.

\subsection{Defining Security Contracts}
A security contract essentially defines what architectural state of the victim is considered exposed or public.
The contract is defined using two building blocks: the observer mode, or what architectural state is exposed, and the execution mode or which paths of a given program are considered (speculatively) reachable and exposed.
We define an observer mode that views the sequence of addresses of every (speculative) victim access to shared memory, and the data of non-speculative accesses, but no other victim architectural state (i.e., the enclave PC is not exposed).
We also consider that the sequence does not contain any timing information (i.e., no timestamps, only data and addresses).
We call this observation mode ``shared memory'' abbreviated to \textbf{\texttt{shm}}.

We now describe the two different execution modes we consider for our contracts.
The first execution mode we consider is the \emph{sequential} mode (no speculation at all, i.e., one-instruction-at-a-time and in-order) or \textbf{\texttt{seq}}, which means that only non-speculative execution paths are considered reachable and addresses and data of shared memory accesses are exposed. 
We call this first contract \textbf{\texttt{seq\_shm}}.
Our second execution mode is the \emph{straight-line speculation} mode or \textbf{\texttt{str}}, which means that we extend the sequential mode to also consider instructions up to a certain depth after a change in control flow (e.g., a branch or a jump) to expose their corresponding observations.
We call this second contract \textbf{\texttt{str\_shm}}.

\subsection{Hardware Side}
We now want to show that each of our hardware mechanisms satisfy one of these contracts, i.e., that they do not leak more information to an attacker that what is stated in the contract.
In our case, we start from a design that satisfies Strong Microarchitectural Isolation.
We also make sure to eliminate on-core microarchitectural state that could entangle state related to enclave private memory, shared memory and the attacker (see Section \ref{sec:l1bypass} where we bypass the L1-D cache).
This means that the attacker is only able to observe the microarchitecture we explicitly share, i.e., the shared LLC partitions.
As a result, the only architectural states the attacker is able to reconstruct through microarchitectural observations are the addresses of speculative shared memory accesses.
Note that per our threat model described in Section \ref{sec:threatmodel} and following previous work, we artificially make the assumption that the attacker is not able to collect timestamps for the accesses it observes.
For our first mechanism, SpecSafe (see Section \ref{sec:mechanism1}), the program is executed following the out-of-order semantics of our processor, but 
we make sure to not touch any shared microarchitecture until a shared memory access has reached the head of the ROB and so is non-speculative.
As a result, Citadel with SpecSafe verifies \textbf{\texttt{seq\_shm}}.
Our second mechanism is the burst mode (see Section \ref{sec:mechanism2}).
Here, the hardware was precisely designed to implement (see Section \ref{sec:mechanism2}) straight-line speculation, and as a result obeys \textbf{\texttt{str\_shm}}.

\noindent 
\begin{minipage}[tbh]{.23\textwidth} % Adjust width as necessary
\begin{lstlisting}[language={[RISC-V]Assembler}, caption={Pathological Insecure Code for Burst Mode}, label={lst:pathologic}]
jal far_away_function
lw a0, 0(secret_address)
\end{lstlisting}
\end{minipage}%
\hfill % Adds horizontal space between minipages
\begin{minipage}[tbh]{.23\textwidth} % Adjust width as necessary
\begin{lstlisting}[language={[RISC-V]Assembler}, caption={Assembly Snippet for the SHA256 library}, label={lst:codesnippet}]
  csrwi 0x802,BURST
.loop1:
[...] # Straight-line code with
        shared memory accesses
  bne a3,a6,.loop1
  csrwi 0x802,NO_BURST
\end{lstlisting}
\end{minipage}

\subsection{Software Side}

    
We now need to analyze if our programs can enforce SNI with regard to these contracts.

In our case, SNI reduces to the fact that the program does not expose more secret data, when considered under a given security contract than when run sequentially, i.e., under \textbf{\texttt{seq\_shm}}.
This is straightforward for our first contract, as the contract is exactly \textbf{\texttt{seq\_shm}}.
Putting everything together, this means that any enclave program secure when run sequentially (i.e., that does not explicitly leak its secrets) is also secure when run in Citadel with SpecSafe.

This is not true for all program for our second contract \textbf{\texttt{str\_shm}}. 
Indeed, we can craft a pathological code snippet (see Listing \ref{lst:pathologic}) that would systematically leak a secret under \textbf{\texttt{str\_shm}} while secure when run sequentially.
Hence, burst mode can be harmful, and it is the responsibility of the programmer to make sure that they use burst mode in a safe context, by analyzing the code under the simple straight-line speculation model.
Fortunately, we only use burst mode on small code snippets and are able to carry out the required program analysis by hand.
This is possible as burst mode is enabled and disabled by writing to a special CSR with a \texttt{csrwi} instruction.
This instruction \textit{cannot} be performed speculatively and acts as a speculation barrier, hence it can be used to constrain the speculative flow control within clear bounds.
Similar to previous work \cite{guarnieri2021hardware, guarnieri2020spectector} we also make the assumption that when entering burst mode, the initial state of the microarchitecture touched by the code is independent of any secrets.
This is reasonable as \texttt{csrwi} already flushes the processor pipeline, predictors are deactivated and the TLB and caches are easy to selectively warm up.
The rest is left to the responsibility of the code before entering Burst mode.

\noindent\textbf{SHA256 Example}
We annotate a \texttt{LOAD} macro in a SHA256 library and show the corresponding assembly snippet in Listing \ref{lst:codesnippet} (full assembly in supplemental material).
The performance impact of this simple annotation can be seen in Section \ref{sec:evalsharedmemory}.
Here, the assembly shows us that the code is enclosed between two \texttt{csrwi}, which bound speculation.
Furthermore, one branch is present, but is immediately fenced by a \texttt{csrwi}. 
It is fenced as \texttt{csrwi} is a system instruction that cannot enter the reorder buffer until the reorder buffer is empty.
As a result, for that specific code snippet, the speculative control flow, and as a result, the exposed architectural states are the same for \textbf{\texttt{str\_shm}} as for \textbf{\texttt{seq\_shm}}, which is sufficient to say that the code snippet realizes SNI with regard to \textbf{\texttt{str\_shm}}.

\noindent\textbf{\texttt{memcpy} Example}
We minimally annotate an implementation of \texttt{memcpy} to take advantage of burst mode when copying data from shared memory to enclave private memory (full C code in in supplemental material).
More complex optimizations could use loop unrolling and annotate the second while loop.
Performance results can be seen in Section \ref{sec:evalsharedmemory}.
We show the details of the assembly for the code snippet we annotate in Figure \ref{fig:memcpyasm}.
Here, the program analysis to verify SNI under \textbf{\texttt{str\_shm}} is more complex as some shared memory accesses are performed speculatively.
Nevertheless, by our definition of SNI, we only care about programs and inputs that lead to a secure execution (i.e., that does not leak secrets) under the \textbf{\texttt{seq\_shm}} contract.
That means that we can assume that registers \texttt{a4} and \texttt{a5} never hold any secret values under \textbf{\texttt{seq\_shm}} as they are exposed to the attacker per our contract, through the shared memory accesses.
Furthermore, at most one more run of \textbf{\texttt{.loop2}} is executed speculatively.
Hence, the only new values an attacker can learn are \texttt{\$a4+8} and \texttt{\$a5+8} with \texttt{\$a4} and \texttt{\$a5} the last values taken by the two registers under \textbf{\texttt{seq\_shm}} and these values are not secret.
Finally, which of these two values are exposed is public as we assume the initial microarchitectural state touched by the program to be public and no secrets are touched.

\noindent\textbf{Program Analysis Tools}
Finally, for more complex pieces of code, static analysis could
be done manually or using many existing tools: e.g.. oo7 \cite{wang2019oo7}, Kasper \cite{johannesmeyer2022kasper}, Kleespectre \cite{wang2020kleespectre}, Spectector \cite{guarnieri2020spectector}.
Note that the usual critique and limitations of these tools are that 1) most code bases are too big to analyze without producing a lot of false positives 2) they make simplifying assumptions on the speculation model.
Because Citadel also leverages hardware isolation, both of those limitations are not concerns. The code snippets we analyze are small and the observation and execution modes are simple, well-defined, and enforced by the hardware.
Hence, these works are well-suited to our needs.
We leave the \textit{automated} analysis of programs by adapting these tools in the context of Citadel as future work (i.e., to support RISCV assembly).